\section{特征标}

\subsection{特征标}

\begin{definition}[特征标]\label{def:character}
设 $G$ 为群, $V$ 为 $\mathbb{F}$-线性空间, $\rho\colon G\to\mathrm{GL}_V$ 为表示.
称 $G$ 上的 $\mathbb{F}$-值函数
\[
  \chi\colon G\to\mathbb{F},\qquad g\longmapsto \operatorname{Tr}(\rho(g))
\]
为表示 $\rho$ 的特征标. 为强调所对应的表示, 常将 $\chi$ 记为 $\chi_\rho$ 或 $\chi_V$.
\end{definition}

\begin{remark}
$\chi_\rho$ 一般不是群同态: 若它是, 就等于要求
$\operatorname{Tr}(AB)=\operatorname{Tr}(A)\operatorname{Tr}(B)$, 而这不成立.
\end{remark}

\begin{lemma}\label{lem:char-class}
设 $V$ 为 $G$ 的表示, 则 $\forall g,h\in G$, $\chi_V(ghg^{-1})=\chi_V(h)$.
\end{lemma}

\begin{proof}
$\chi_V(ghg^{-1})=\operatorname{Tr}\bigl(\rho(ghg^{-1})\bigr)
=\operatorname{Tr}\bigl(\rho(g)\rho(h)\rho(g)^{-1}\bigr)=\operatorname{Tr}\bigl(\rho(h)\bigr)=\chi_V(h)$.
\end{proof}

\begin{definition}[类函数]\label{def:class-function}
称函数 $f\colon G\to\mathbb{F}$ 为 $G$ 上的一个类函数, 若 $f$ 在一个共轭类上取常值,
即在同一共轭类上取值不变.
\end{definition}

\begin{remark}
由 Lemma \ref{lem:char-class}, $\chi_V$ 为类函数.
\end{remark}

\begin{lemma}\label{lem:char-inverse}
设 $\mathbb{F}=\mathbb{C}$, $\rho\colon G\to\mathrm{GL}_V$ 为表示, 则
$\chi_\rho(g^{-1})=\overline{\chi_\rho(g)}$.
\end{lemma}

\begin{proof}
设 $g^m=1$, 则 $\rho(g)^m=\rho(g^m)=\mathrm{Id}$. 于是 $\rho(g)$ 可对角化, 其本征值
$\lambda_1,\dots,\lambda_m$ 满足 $\lambda_i^m=1$, 从而 $\lambda_i^{-1}=\overline{\lambda_i}$. 故
\[
  \chi_\rho(g^{-1})=\operatorname{Tr}\bigl(\rho(g)^{-1}\bigr)=\sum_i\lambda_i^{-1}
  =\sum_i\overline{\lambda_i}=\overline{\operatorname{Tr}(\rho(g))}=\overline{\chi_\rho(g)}.
\]
\end{proof}

\subsection{$G$ 上的函数与 $\mathbb{F}G$ 的对偶}

\begin{proposition}\label{prop:dual-is-functions}
$(\mathbb{F}G)^*$ 就是 $G$ 上所有的函数.
\end{proposition}

\begin{proof}
对函数 $f\colon G\to\mathbb{F}$, 把 $f$ 线性延拓为 $\mathbb{F}G\to\mathbb{F}$ 上的线性函数,
即得 $f\in(\mathbb{F}G)^*$; 反之 $(\mathbb{F}G)^*$ 中元素限制在基 $G$ 上即为 $G$ 上的函数.
\end{proof}

\begin{lemma}\label{lem:dual-algebra}
$(\mathbb{F}G)^*$ 是代数.
\end{lemma}

\begin{proof}
$(\mathbb{F}G)^*$ 自然是线性空间. 定义其上的乘法: $\forall f_1,f_2\in(\mathbb{F}G)^*$,
\[
  (f_1\cdot f_2)(g)=f_1(g)\cdot f_2(g),\qquad g\in G .
\]
该乘法有单位元 $\mathbf{1}\colon G\to\mathbb{F}$, $g\mapsto 1$. 结合律与双线性性逐点验证即可.
\end{proof}

\begin{definition}[$*$-代数]\label{def:star-algebra}
称代数 $A$ 为 $*$-代数, 若存在线性变换 $*\colon A\to A$ 使得
\begin{enumerate}
  \item $*\circ *=\mathrm{Id}$;
  \item $\forall a,b\in A$, $*(ab)=*(b)\cdot *(a)$.
\end{enumerate}
\end{definition}

\begin{example}
全矩阵代数 $M_n(\mathbb{F})$ 在转置下是 $*$-代数:
\[
  M_n(\mathbb{F})\to M_n(\mathbb{F}),\qquad A\longmapsto A^{\mathrm{T}},
\]
有 $(A^{\mathrm{T}})^{\mathrm{T}}=A$ 且 $(AB)^{\mathrm{T}}=B^{\mathrm{T}}A^{\mathrm{T}}$.
\end{example}

\begin{proposition}\label{prop:dual-star}
$(\mathbb{F}G)^*$ 为 $*$-代数.
\end{proposition}

\begin{proof}
定义 $*\colon(\mathbb{F}G)^*\to(\mathbb{F}G)^*$, $f\mapsto *(f)$, 其中
$\bigl(*(f)\bigr)(g)=f(g^{-1})$. 则
$\bigl(*(*(f))\bigr)(g)=\bigl(*(f)\bigr)(g^{-1})=f(g)$, 即 $*\circ *=\mathrm{Id}$.
其次对 $f_1,f_2\in(\mathbb{F}G)^*$,
\[
  \bigl(*(f_1f_2)\bigr)(g)=(f_1f_2)(g^{-1})=f_1(g^{-1})f_2(g^{-1})
  =\bigl(*(f_2)\bigr)(g)\bigl(*(f_1)\bigr)(g),
\]
即 $*(f_1f_2)=*(f_2)\cdot *(f_1)$.
\end{proof}

\begin{remark}
回忆 $V$ 上的双线性函数 $(\cdot,\cdot)\colon V\times V\to\mathbb{F}$.
称其是对称的, 若 $(\alpha,\beta)=(\beta,\alpha)$, $\alpha,\beta\in V$;
称其是非退化的, 若 $\forall\beta\in V$, $(\alpha,\beta)=0$ 蕴含 $\alpha=0$.
$\mathbb{R}^n$ 上的内积即为对称非退化的双线性函数.
\end{remark}

\begin{proposition}\label{prop:dual-pairing}
设 $\operatorname{char}\mathbb{F}\nmid |G|$, 则 $(\mathbb{F}G)^*$ 上存在对称非退化的双线性函数.
\end{proposition}

\begin{proof}
定义
\[
  (\cdot,\cdot)\colon(\mathbb{F}G)^*\times(\mathbb{F}G)^*\to\mathbb{F},\qquad
  (f_1,f_2)=\frac{1}{|G|}\sum_{g\in G}f_1(g)f_2(g^{-1}).
\]
双线性性由逐点运算直接得到. 对称性:
\[
  (f_1,f_2)=\frac{1}{|G|}\sum_{g}f_1(g)f_2(g^{-1})
  =\frac{1}{|G|}\sum_{g}f_2(g)f_1(g^{-1})=(f_2,f_1).
\]
非退化性: 设 $f_1\neq 0$, 取 $g\in G$ 使 $f_1(g)\neq 0$, 并令 $\delta_g\in(\mathbb{F}G)^*$
为 $\delta_g(g')=1$ ($g'=g$), $\delta_g(g')=0$ ($g'\neq g$). 则
\[
  (f_1,\delta_g)=\frac{1}{|G|}\sum_{h}f_1(h)\delta_g(h^{-1})=\frac{1}{|G|}f_1(g)\neq 0 .
\]
故 $(\cdot,\cdot)$ 非退化.
\end{proof}

\subsection{特征标的正交性}

先补充一个记号:
\[
  \operatorname{Hom}_G(V_1,V_2)=\bigl\{\varphi\colon V_1\to V_2 \bigm| \varphi \text{ 线性, 且 }
  \varphi(g\cdot v_1)=g\cdot\varphi(v_1),\ g\in G,\ v_1\in V_1\bigr\}.
\]
下面以几个特征标的计算引入.

\begin{example}\label{eg:char-Z3}
(1) 回忆 $\mathbb{C}\mathbb{Z}_3$. 在 II.\ 例 2 中我们讨论了 $\mathbb{C}\mathbb{Z}_n$, 利用结论可写
\[
  \mathbb{C}\mathbb{Z}_3=\bigoplus_{j=1}^{3}\mathbb{C}x_j,\qquad
  \mathbb{C}x_j=\operatorname{span}_{\mathbb{C}}\{x_j\},
\]
其中
\[
  x_1=\frac{1}{3}(1+\omega^2g+\omega g^2),\quad
  x_2=\frac{1}{3}(1+\omega g+\omega^2g^2),\quad
  x_3=\frac{1}{3}(1+g+g^2).
\]
记 $\chi_{\mathbb{C}x_2}\colon\mathbb{Z}_3\to\mathbb{C}$ 为
$1\mapsto 1$, $g\mapsto\omega^2$, $g^2\mapsto\omega^4=\omega$;
$\chi_{\mathbb{C}x_1}\colon\mathbb{Z}_3\to\mathbb{C}$ 为
$1\mapsto 1$, $g\mapsto\omega$, $g^2\mapsto\omega^2$. 则
\[
  (\chi_{\mathbb{C}x_2},\chi_{\mathbb{C}x_1})
  =\frac{1}{3}\sum_{g\in\mathbb{Z}_3}\chi_{\mathbb{C}x_2}(g)\chi_{\mathbb{C}x_1}(g^{-1})
  =\frac{1}{3}(1+\omega+\omega^2)=0,
\]
而注意 $(\chi_{\mathbb{C}x_2},\chi_{\mathbb{C}x_2})=(\chi_{\mathbb{C}x_1},\chi_{\mathbb{C}x_1})=1$.

(2) 回忆 $S_3=\{(),(12),(13),(23),(123),(132)\}$,
\[
  \rho_{\mathrm{triv}}\colon S_3\to\mathbb{C}^\times,\ g\mapsto 1;\qquad
  \rho_{\mathrm{sign}}\colon S_3\to\mathbb{C}^\times,\
  \text{奇置换}\mapsto -1,\ \text{偶置换}\mapsto 1 .
\]
再看 $\rho_{\mathrm{std}}\colon S_3\to\mathrm{GL}_2$: 标准表示作用在
\[
  V_1=\Bigl\{\sum_i k_ie_i \Bigm| \sum_i k_i=0\Bigr\}
\]
上, 取 $V_1$ 的一组基 $\{e_1-e_2,\ e_1-e_3\}$. 显然 $()\mapsto I_2$, 且
\[
  (12)(e_1-e_2,\ e_1-e_3)=(e_2-e_1,\ e_2-e_3)
  =(e_1-e_2,\ e_1-e_3)\begin{pmatrix}-1&-1\\ 0&1\end{pmatrix},
\]
故 $(12)\mapsto\begin{pmatrix}-1&-1\\ 0&1\end{pmatrix}$; 而 $(13),(23)$ 与 $(12)$ 共轭.
由特征标是类函数, 用 $(12)$ 计算的结果与用 $(13),(23)$ 计算的一致, 此处从略.
\end{example}

类似地, 计算 $(123)$:
\[
  (123)(e_1-e_2,\ e_1-e_3)=(e_2-e_3,\ e_2-e_1)
  =(e_1-e_2,\ e_1-e_3)\begin{pmatrix}-1&-1\\ 1&0\end{pmatrix}.
\]
现在有了具体的表示, 便可计算特征标. 由上面的讨论,
\[
  (\chi_{\mathrm{triv}},\chi_{\mathrm{triv}})
  =\frac{1}{6}\sum_{g\in S_3}\chi_{\mathrm{triv}}(g)\chi_{\mathrm{triv}}(g^{-1})
  =\frac{1}{6}\sum_{g\in S_3}1=1 .
\]
注意到 $S_3$ 中 $(),(123),(132)$ 是偶置换, $(12),(13),(23)$ 是奇置换, 于是
\[
  (\chi_{\mathrm{triv}},\chi_{\mathrm{sign}})
  =\frac{1}{6}\sum_{g}\chi_{\mathrm{triv}}(g)\chi_{\mathrm{sign}}(g^{-1})
  =\frac{1}{6}\bigl(1-1-1-1+1+1\bigr)=0,
\]
\[
  (\chi_{\mathrm{triv}},\chi_{\mathrm{std}})
  =\frac{1}{6}\sum_{g}\chi_{\mathrm{triv}}(g)\chi_{\mathrm{std}}(g^{-1})
  =\frac{1}{6}\bigl(2+0+0+0-1-1\bigr)=0 .
\]
类似可得
\[
  (\chi_{\mathrm{sign}},\chi_{\mathrm{sign}})=1,\qquad
  (\chi_{\mathrm{sign}},\chi_{\mathrm{std}})=0,\qquad
  (\chi_{\mathrm{std}},\chi_{\mathrm{std}})=1 .
\]
似乎出现了某种正交性, 下面逐步说明.

\begin{lemma}\label{lem:averaging}
设 $G$ 为群, $\rho\colon G\to\mathrm{GL}_V$ 为表示, 记
$\zeta=\dfrac{1}{|G|}\sum_{g\in G}\rho(g)$, 则 $\zeta\in\operatorname{Hom}_G(V,V)$.
\end{lemma}

\begin{proof}
只需证明 $\zeta(g\cdot v)=g\cdot\zeta(v)$, $g\in G$. 事实上
\[
  \zeta(g\cdot v)=\frac{1}{|G|}\sum_{h}\rho(h)\rho(g)(v)
  =\frac{1}{|G|}\sum_{h}\rho(hg)(v)=\zeta(v),
\]
同理 $g\cdot\zeta(v)=\dfrac{1}{|G|}\sum_{h}\rho(g)\rho(h)(v)
=\dfrac{1}{|G|}\sum_{h}\rho(gh)(v)=\zeta(v)$. 故 $\zeta(g\cdot v)=g\cdot\zeta(v)$.
\end{proof}

\begin{corollary}\label{cor:zeta-commute}
$\rho(g)\cdot\zeta=\zeta\cdot\rho(g)=\zeta$.
\end{corollary}

\begin{definition}\label{def:invariants}
设 $G$ 为群, $V$ 为 $G$ 的表示, 定义
\[
  V^G=\{v\in V\mid g\cdot v=v,\ \forall g\in G\},
\]
这是 $V$ 的一个子表示.
\end{definition}

\begin{remark}
$V^G$ 是 $V$ 的平凡子表示的和 (而 $V^G$ 本身也是平凡的).
\end{remark}

\begin{lemma}\label{lem:zeta-idempotent}
$\zeta^2=\zeta$.
\end{lemma}

\begin{proof}
由 $\zeta=\dfrac{1}{|G|}\sum_{g}\rho(g)$ 及 Corollary \ref{cor:zeta-commute},
$\rho(g)\zeta=\zeta$, 故
\[
  \zeta^2=\frac{1}{|G|}\sum_{g}\rho(g)\zeta=\frac{1}{|G|}\sum_{g}\zeta
  =\frac{1}{|G|}\cdot|G|\cdot\zeta=\zeta .
\]
\end{proof}

\begin{lemma}\label{lem:zeta-image}
$\zeta(V)=V^G$.
\end{lemma}

\begin{proof}
(1) 任取 $v\in V$, 则 $\zeta(v)\in V$. 对任意 $g\in G$ 有 $g\cdot\zeta(v)=\zeta(v)$,
故 $\zeta(v)\in V^G$, 即 $\zeta(V)\subseteq V^G$.

(2) 取 $v\in V^G$, 则
$\zeta(v)=\dfrac{1}{|G|}\sum_{g}\rho(g)(v)=\dfrac{1}{|G|}\sum_{g}v=v$,
即 $v=\zeta(v)\in\zeta(V)$, 故 $V^G\subseteq\zeta(V)$.

综上 $\zeta(V)=V^G$.
\end{proof}

\begin{proposition}\label{prop:dim-invariants}
设 $G$ 为群, $\rho\colon G\to\mathrm{GL}_V$ 为表示, 则
\[
  \frac{1}{|G|}\sum_{g\in G}\chi_V(g)
  =\frac{1}{|G|}\sum_{g\in G}\chi_{V^G}(g)=\dim V^G .
\]
\end{proposition}

\begin{proof}
由 Lemma \ref{lem:zeta-image}, $\zeta\colon V\to V^G$; 由 Lemma \ref{lem:zeta-idempotent}
及 Lemma \ref{lem:averaging}, $\zeta|_{V^G}=\mathrm{Id}_{V^G}$.
于是 $V=\operatorname{Ker}(\zeta)\oplus V^G$, 从而 $\operatorname{Tr}(\zeta)=\operatorname{Tr}(\zeta|_{V^G})$. 故
\[
  \frac{1}{|G|}\sum_{g}\chi_V(g)=\operatorname{Tr}(\zeta)
  =\operatorname{Tr}(\zeta|_{V^G})=\frac{1}{|G|}\sum_{g}\chi_{V^G}(g)
  =\frac{1}{|G|}\sum_{g}\dim V^G=\dim V^G .
\]
\end{proof}

\begin{theorem}\label{thm:inner-product-dim}
设 $V_1,V_2$ 为 $G$ 的两个表示, 则
$(\chi_{V_1},\chi_{V_2})=\dim\operatorname{Hom}_G(V_1,V_2)$.
\end{theorem}

\begin{proof}
由 Lemma \ref{lem:char-inverse} 与 Proposition \ref{prop:dim-invariants},
\[
  (\chi_{V_1},\chi_{V_2})=\frac{1}{|G|}\sum_{g}\chi_{V_1}(g^{-1})\chi_{V_2}(g)
  =\frac{1}{|G|}\sum_{g}\chi_{V_1^*\otimes V_2}(g)=\dim\,(V_1^*\otimes V_2)^G .
\]
注意 $V_1^*\otimes V_2\cong\operatorname{Hom}(V_1,V_2)$, 故结论成立.
\end{proof}
